\chapter*{Abstract}
\addcontentsline{toc}{chapter}{Abstract}
Predicting concrete compressive strength from mixture composition and curing age is a nonlinear regression problem. Support Vector Regression (SVR) can model these relationships, but its performance depends on the regularization parameter, radial basis function kernel width and epsilon-insensitive loss threshold. This project evaluates whether Differential Evolution (DE) provides a useful automated method for selecting these hyperparameters under a fixed computational budget.

The study uses the UCI Concrete Compressive Strength dataset with 1,030 observations, eight numerical inputs and compressive strength as the target in megapascals. Identical input tuples are assigned to common groups before splitting. The fixed partition contains 826 development observations and 204 holdout observations. A StandardScaler--SVR pipeline is fitted separately inside each of five group-disjoint development folds, avoiding leakage from preprocessing or repeated input combinations. DE/rand/1/bin is implemented explicitly with differential mutation, binomial crossover, forced mutant-coordinate inheritance, clipped bounds and deferred greedy selection.

Default RBF-SVR is compared with Random Search and DE using the same three-dimensional logarithmic search space. Each tuned method receives 300 candidate evaluations per seed, including initialization, for five predefined optimizer seeds. Candidate selection uses mean development cross-validation RMSE. The original baseline and Random Search holdout evaluations preceded final DE evaluation; this chronology is retained rather than describing the entire study as prospectively blinded.

Default SVR obtains holdout RMSE of 9.9069 MPa, MAE of 7.8854 MPa and an $R^2$ of 0.5755. Random Search achieves mean RMSE of $6.2328\pm0.1194$ MPa, while DE achieves $6.2250\pm0.1104$ MPa across optimizer seeds. DE reduces mean holdout RMSE by approximately 37.2\% relative to default SVR. Its difference from Random Search is only about 0.008 MPa, so the results support the value of hyperparameter tuning without establishing significant superiority of DE. Convergence, prediction-residual and development learning-curve diagnostics accompany the comparison. The contribution is a reproducible, budget-controlled experimental pipeline and an auditable implementation of DE; broader deployment and statistical generalization require further validation.

\noindent\textbf{Keywords:} concrete compressive strength; Differential Evolution; Support Vector Regression; hyperparameter optimization; group cross-validation; reproducibility.
